% !PRINTED 270
\weberpage{270}
\begin{center}
\gesperrt{Achter Abschnitt.}

\gesperrt{Der Sturm'sche Lehrsatz.}

\rule{0.3\linewidth}{0.4pt}

\S.\,85.

\gesperrt{Das Sturm'sche Problem.}
\end{center}

Die im vorigen Abschnitt behandelte Frage nach der Anzahl der reellen Wurzeln einer Gleichung ist ein specieller Fall eines allgemeineren Problems, das den Gegenstand dieses Abschnittes ausmachen soll, und das die Grundlage ist für alle Methoden der genäherten numerischen Berechnung von Gleichungswurzeln.

Es handelt sich um die Frage: \gesperrt{wie viele reelle Wurzeln einer reellen numerischen Gleichung liegen zwischen zwei gegebenen reellen Zahlwerthen} $a$ und $b$?

Ist dies entschieden, so handelt es sich weiter darum, die Grenzen $a$, $b$ so weit einzuengen, dass nur noch \gesperrt{eine Wurzel} der gegebenen Gleichung zwischen ihnen liegt, und sie endlich einander so weit zu nähern, dass jede von ihnen als ein \gesperrt{genäherter Werth} dieser Wurzel betrachtet werden kann.

Nehmen wir $a = -\infty$, $b = +\infty$, oder doch $a$ negativ, $b$ positiv so gross an, dass jenseits dieser Grenzen keine Wurzeln mehr liegen können, so fällt diese Aufgabe mit der im vorigen Abschnitt behandelten, die Anzahl aller reellen Wurzeln zu bestimmen, zusammen.

Wir wollen zunächst zeigen, wie sich das allgemeine Problem auf das specielle zurückführen lässt, wie also unsere jetzt aufgeworfene Frage, im Princip wenigstens, durch die Betrachtungen des vorigen Abschnittes beantwortet ist. Es möge sich zunächst
%CONT>>
